\chapter{Consejos y avisos de seguridad}
\label{cap:seguridad}

\chapterquote{Estudiar los errores de otros pilotos le ayudará a no cometerlos.}{Robinson Helicopter Company}

\section{Generalidades}
\fuente{POH, sección 10, y Safety Notices de Robinson revisadas hasta julio de 2025}

Este capítulo entrega sugerencias diversas para operar el helicóptero con mayor seguridad, el programa de conocimientos y competencia recomendado y el texto completo de las \emph{Safety Notices} (avisos de seguridad, SN) emitidas por Robinson a raíz de accidentes e incidentes.

\section{Consejos de seguridad}

\begin{pasos}
\item \textbf{Nunca empuje el cíclico hacia adelante para descender o para terminar un tirón hacia arriba} (como se haría en un avión). Puede producir una condición de baja G (casi ingravidez) que puede hacer que una pala del rotor principal golpee la cabina. Inicie siempre el descenso con el colectivo (SN-11).
\item Nunca deje intencionalmente que el combustible baje tanto en vuelo que se encienda la luz de bajo combustible (SN-15).
\item Nunca deje el helicóptero sin protección donde curiosos puedan dañar inadvertidamente partes críticas, como las palas del rotor de cola.
\item Encienda la luz estroboscópica antes de arrancar el motor y déjela encendida hasta que los rotores se detengan. Está cerca del rotor de cola y advierte al personal de tierra. Se recomienda operar todas las luces exteriores en vuelo, incluso de día, para favorecer la prevención de colisiones.
\item Nunca lleve carga externa salvo con un gancho aprobado, y nunca fije nada al exterior del helicóptero. Asegúrese de que no haya objetos sueltos en la cabina, sobre todo al volar sin puertas: incluso un objeto pequeño y liviano puede dañar el rotor de cola en vuelo.
\item Evite acciones bruscas de los controles y maniobras aceleradas, sobre todo a alta velocidad. Producen altas cargas de fatiga que, con el tiempo, podrían llevar a la falla de un componente crítico.
\item Un cambio en el sonido o la vibración del helicóptero puede indicar la falla inminente de un componente crítico. Si en vuelo comienza un sonido o vibración inusual, haga un aterrizaje seguro y haga inspeccionar a fondo la aeronave antes de reanudar el vuelo. Haga un estacionario cerca del suelo para verificar que el problema se resolvió y luego haga reinspeccionar la aeronave antes de reanudar el vuelo libre.
\item Asegúrese de que el personal de tierra y los curiosos se mantengan bien alejados del rotor de cola y del escape. Las palas del rotor principal también pueden ser peligrosas, especialmente si hay personas en terreno más alto que el helicóptero.
\item Nunca deje que la RPM del rotor baje peligrosamente. La mayoría de los aterrizajes duros son sobrevivibles mientras no se deje entrar el rotor en pérdida.
\item Nunca despegue ni aterrice con viento de cola, especialmente en altura. La pérdida de sustentación traslacional resultante puede hacer que el helicóptero se hunda contra el suelo o los obstáculos.
\item Un descenso vertical o una aproximación empinada, sobre todo con viento de cola, puede hacer que el rotor vuele dentro de su propio flujo descendente: el estado de anillo de vórtice. En esta condición, incluso la potencia máxima del motor puede no bastar para detener el descenso, con un impacto duro contra el suelo. Se evita reduciendo siempre la razón de descenso antes de reducir la velocidad (SN-22).
\item El helicóptero es estable sobre su tren mientras el contacto con el suelo sea vertical o con el helicóptero avanzando. Si toca moviéndose hacia atrás o de lado, pueden producirse daños e incluso un vuelco. Los pilotos con poca experiencia y los alumnos deben practicar los aterrizajes y el estacionario con el helicóptero avanzando lentamente.
\item Al operar en altura (sobre 3000 o 4000~ft), el acelerador suele estar completamente abierto y la RPM debe controlarse con el colectivo. La correlación acelerador/colectivo no es eficaz en estas condiciones y la respuesta del gobernador es bastante lenta. Puede ser necesario cerrar el acelerador al bajar el colectivo para evitar una sobrevelocidad.
\item No use el paso colectivo para frenar el rotor durante la detención: genera sustentación en las palas, que puede vencer la fricción de los topes de caída y permitir que las palas golpeen el cono de cola. Tampoco frene ni detenga los rotores tomando el rotor de cola con la mano: puede dañar la transmisión del rotor de cola.
\item No aterrice en pasto alto y seco. El escape está cerca del suelo y muy caliente; puede iniciarse un incendio.
\item Revise siempre el área en busca de cables u otras obstrucciones antes de practicar autorrotaciones.
\item Nunca haga una prueba de magnetos en vuelo. Si un magneto funciona mal y el piloto pone a masa el otro, el motor puede detenerse por completo. Si sospecha una falla de magneto, deje la llave en \textsf{BOTH} y aterrice tan pronto como sea práctico.
\end{pasos}

\section{Conocimiento y competencia del piloto}

El conocimiento y la competencia del piloto son esenciales para operar el helicóptero con seguridad. Además de tener las licencias correspondientes y cumplir los requisitos reglamentarios de recurrencia (como el 14 CFR \S\S\,61.56 y 61.57), todos los pilotos deben mantener actualizados sus conocimientos y su competencia a un alto nivel.

Se necesitan estudio periódico y entrenamiento recurrente. Se recomienda poner énfasis en las áreas siguientes, que también deben cubrirse en el entrenamiento de transición a cada marca y modelo. Todos los distribuidores Robinson cuentan con instructores aprobados, y la fábrica puede informar sobre instructores en su zona.

\begin{nota}
Los \emph{Helicopter Practical Test Standards} y el \emph{Helicopter Flying Handbook} de la FAA, así como los programas comerciales de instrucción en helicópteros, entregan lecciones y estándares detallados para muchas de estas áreas.
\end{nota}

\begin{itemize}
\item Todos los consejos de seguridad y Safety Notices del POH.
\item Limitaciones y procedimientos de emergencia.
\item Estacionario de precisión, giros en estacionario y rodaje en estacionario (SN-9 y SN-42): despegue y posado seguros sin tendencia a derivar; control de posición y de guiñada con viento cruzado y de cola.
\item Estado de anillo de vórtice (SN-22 y SN-34): condiciones propicias; procedimientos de recuperación (Vuichard y tradicional).
\item Pérdida de referencias visuales exteriores (SN-18, SN-19 y SN-26): la gravedad de esta condición; las trampas (vuelo nocturno con nubes, condiciones que empeoran gradualmente, etc.); terreno sin rasgos o agua espejada.
\item Planificación del vuelo (SN-15, SN-26 y SN-43): inspección prevuelo minuciosa; combustible; meteorología; performance (calor, altura, carga).
\item Distracciones (SN-16, SN-34, SN-36 y SN-41): no mantener la vista afuera buscando cables, otros obstáculos y tránsito; misiones de alta carga de trabajo como vuelos fotográficos; pasajeros; aviónica; teléfonos celulares.
\item Baja G y golpe de mástil (SN-11, SN-29 y SN-32): cómo evitarlos; reducir la velocidad en turbulencia; vigilar la velocidad con poca carga; asegurarse de que se hayan retirado los controles del pasajero; reconocimiento y recuperación.
\end{itemize}

\begin{precaucion}
Nunca practique ni demuestre la baja G en vuelo. El entrenamiento de baja G debe ser solo teórico.
\end{precaucion}

\begin{itemize}
\item Consideraciones de baja RPM (SN-10, SN-24 y SN-29): reconocimiento y recuperación.
\item Fallas de potencia (SN-10, SN-24 y SN-29): entrada instintiva en autorrotación; considerar continuamente lugares de aterrizaje de emergencia en todo vuelo.
\item Autorrotaciones de práctica (SN-38): métodos probados y seguros.
\end{itemize}

\begin{precaucion}
La práctica en vuelo de baja RPM, fallas de potencia y autorrotaciones solo debe hacerse bajo la supervisión de un instructor.
\end{precaucion}

\begin{itemize}
\item Hielo en el carburador (SN-25 y SN-31): condiciones propicias; uso de la calefacción.
\item Operaciones en áreas confinadas (SN-22): reconocimiento alto y bajo; evaluación del viento; márgenes de potencia.
\end{itemize}

\section{Índice de Safety Notices}

\begin{longtable}{@{}lp{\dimexpr\linewidth-1.4cm-2\tabcolsep}@{}}
\toprule
SN & Título \\
\midrule
\endhead
SN-1 & Accionamiento inadvertido del control de mezcla en vuelo \\
SN-9 & Muchos accidentes involucran vuelco dinámico \\
SN-10 & Accidentes fatales causados por pérdida del rotor por baja RPM \\
SN-11 & Empujes en baja G: extremadamente peligrosos \\
SN-13 & No fije objetos a los patines \\
SN-15 & Quedarse sin combustible puede ser fatal \\
SN-16 & Los tendidos eléctricos son mortales \\
SN-17 & Nunca abandone el helicóptero con el motor en marcha / Sujete los controles al embarcar pasajeros / Nunca aterrice en pasto alto y seco \\
SN-18 & La pérdida de visibilidad puede ser fatal / El exceso de confianza predomina en los accidentes \\
SN-19 & Volar bajo sobre el agua es muy peligroso \\
SN-20 & Cuidado con los vuelos de demostración o de instrucción inicial \\
SN-22 & Estado de anillo de vórtice: cómo evitarlo, reconocerlo y recuperarse \\
SN-23 & Caminar hacia el rotor de cola puede ser fatal \\
SN-24 & La pérdida del rotor por baja RPM puede ser fatal \\
SN-25 & Hielo en el carburador \\
SN-26 & Vuelo nocturno con mal tiempo puede ser mortal \\
SN-27 & Los cortes de acelerador por sorpresa pueden ser mortales \\
SN-28 & Escuche el aviso de una falla inminente de rodamiento / Advertencia de la luz del embrague \\
SN-29 & Pilotos de avión: alto riesgo al volar helicópteros \\
SN-30 & Los objetos sueltos pueden ser fatales \\
SN-31 & El gobernador puede enmascarar el hielo en el carburador \\
SN-32 & Vientos fuertes o turbulencia \\
SN-33 & Holgura de las correas de transmisión \\
SN-34 & Vuelos de levantamiento aéreo y fotográficos: riesgo muy alto \\
SN-35 & Vuelo cerca de torres de radiodifusión \\
SN-36 & Sobrevelocidades durante el despegue \\
SN-37 & Exceder las limitaciones aprobadas puede ser fatal \\
SN-38 & Las autorrotaciones de práctica causan muchos accidentes de instrucción \\
SN-39 & Una vibración inusual puede indicar una grieta en una pala del rotor principal \\
SN-40 & Incendios posteriores al accidente \\
SN-41 & Distracciones del piloto \\
SN-42 & Guiñada imprevista \\
SN-43 & Extreme las precauciones en los vuelos posteriores al mantenimiento \\
SN-44 & Transporte de pasajeros \\
\bottomrule
\end{longtable}

Las SN-2 a SN-8 fueron reemplazadas o eliminadas; la SN-12 fue reemplazada por la SN-24; la SN-14, por las SN-17, SN-27 y SN-28; la SN-21 fue eliminada.

\section{Texto de las Safety Notices}

\subsection*{SN-1 \quad Accionamiento inadvertido del control de mezcla en vuelo}
{\small\textit{Emitida: enero de 1981. Revisada: febrero de 1989; junio de 1994.}}

Se han informado casos en que un piloto tiró inadvertidamente del control de mezcla en lugar de la calefacción del carburador u otro control, lo que produjo la detención súbita y completa del motor. Los pomos tienen formas distintas y el control de mezcla tiene un protector que debe retirarse y un botón de traba que debe presionarse antes de accionarlo. Estas diferencias deben recalcarse al instruir a pilotos nuevos. En el R22 es buena práctica acceder siempre a la compensación lateral pasando la mano por el lado izquierdo del cíclico, lo que reduce la posibilidad de tirar de la mezcla por error. Use siempre el pequeño protector plástico que se coloca en el control de mezcla antes de arrancar el motor y que no se retira hasta el final del vuelo, al cortar la mezcla. Vuelva a colocar el protector para que esté en su lugar en el vuelo siguiente.

Si se tira inadvertidamente de la mezcla, baje el colectivo y entre en autorrotación. Si hay altura suficiente, empuje la mezcla y vuelva a arrancar el motor con la mano izquierda. \textbf{NO} desacople el embrague.

\subsection*{SN-9 \quad Muchos accidentes involucran vuelco dinámico}
{\small\textit{Emitida: julio de 1982. Revisada: junio de 1994.}}

Puede producirse un vuelco dinámico siempre que el tren toque un objeto fijo que obligue al helicóptero a pivotar sobre el objeto en lugar de hacerlo sobre su propio centro de gravedad. El objeto fijo puede ser cualquier obstáculo o superficie que impida que el patín se desplace lateralmente. Una vez iniciado, el vuelco dinámico no puede detenerse solo con cíclico contrario. Por ejemplo, si el patín derecho toca un objeto y se convierte en el punto de pivote mientras el helicóptero empieza a alabear a la derecha, incluso con cíclico izquierdo a fondo el vector de empuje del rotor pasará por la izquierda del punto de pivote y producirá un momento de alabeo hacia la derecha en lugar de hacia la izquierda. El vector y su momento acompañarán al helicóptero mientras siga alabeando a la derecha. \textbf{Bajar rápidamente el colectivo es la forma más eficaz de detener un vuelco dinámico.}

Para evitarlo:
\begin{pasos}
\item Practique siempre las autorrotaciones en estacionario contra el viento, y nunca con viento racheado o superior a 10~nudos.
\item Nunca haga estacionario cerca de cercos, aspersores, arbustos, luces de pista u otros obstáculos donde pueda engancharse un patín.
\item Use siempre un despegue en dos pasos: suba solo el colectivo necesario para quedar liviano sobre los patines y sienta el equilibrio; luego eleve suavemente el helicóptero.
\item No practique maniobras en estacionario cerca del suelo. Mantenga los patines al menos a cinco pies del suelo al practicar vuelo lateral o hacia atrás.
\end{pasos}

\subsection*{SN-10 \quad Accidentes fatales causados por pérdida del rotor por baja RPM}
{\small\textit{Emitida: octubre de 1982. Revisada: febrero de 1989; junio de 1994.}}

Una causa principal de accidentes fatales en helicópteros livianos es no mantener la RPM del rotor. Para evitarlo, todo piloto debe tener sus reflejos condicionados para \textbf{abrir el acelerador y bajar el colectivo instantáneamente} para mantener la RPM ante cualquier emergencia.

El R22 y el R44 han demostrado una excelente resistencia al impacto siempre que el piloto vuele la aeronave hasta el suelo y haga una recogida al final para reducir la velocidad y la razón de descenso. Incluso al descender hacia terreno irregular, árboles, cables o agua, debe obligarse a bajar el colectivo para mantener la RPM hasta justo antes del impacto. El helicóptero puede volcarse y dañarse gravemente, pero los ocupantes tienen una excelente probabilidad de salir ilesos.

La potencia disponible del motor es directamente proporcional a la RPM. Si la RPM cae un 10\,\%, hay un 10\,\% menos de potencia. Con menos potencia el helicóptero empieza a hundirse y, si se sube el colectivo para detenerlo, la RPM bajará aún más y el hundimiento será más rápido. Si el piloto no solo no baja el colectivo sino que lo sube para evitar el descenso, el rotor entrará en pérdida casi de inmediato. En pérdida, las palas se ``soplan hacia atrás'' (\emph{blow back}) y cortan el cono de cola, o simplemente dejan de volar y el helicóptero cae a una velocidad extrema. En ambos casos el accidente probablemente será fatal.

Cualquiera sea la causa de la baja RPM, el piloto debe primero abrir el acelerador y bajar el colectivo simultáneamente para recuperar la RPM antes de investigar el problema. Debe ser un reflejo condicionado. En vuelo de avance, aplicar cíclico atrás para reducir la velocidad también ayuda a recuperar la RPM perdida.

\subsection*{SN-11 \quad Empujes en baja G: extremadamente peligrosos}
{\small\textit{Emitida: octubre de 1982. Revisada: noviembre de 2000; julio de 2025.}}

Empujar el cíclico hacia adelante después de un tirón hacia arriba o un ascenso rápido ---o, en algunas condiciones, incluso desde vuelo nivelado--- puede producir una condición de vuelo en baja G (ingravidez). En baja G el empuje del disco rotor disminuye, y con él la capacidad de generar los momentos necesarios para controlar la actitud. Los helicópteros tienen poca inercia rotacional en el eje de alabeo y cualquier desequilibrio de las fuerzas aerodinámicas producirá un alabeo acelerado. La severidad del alabeo aumenta notablemente con la velocidad y con poca carga, y depende fuertemente del tipo de estabilizador horizontal:

\begin{itemize}
\item Los helicópteros con el \textbf{estabilizador horizontal asimétrico} (versión anterior, montado junto a la caja de cola) pueden experimentar un alabeo rápido a la derecha en baja G.
\item Los helicópteros con el \textbf{estabilizador horizontal simétrico} (versión posterior, montado bajo el cono de cola por delante del rotor de cola) pueden experimentar un alabeo menos severo a la derecha \emph{o} a la izquierda.
\end{itemize}

Con el empuje reducido, la menor eficacia de los controles puede llevar al piloto a hacer grandes acciones para contrarrestar el alabeo, y estas acciones inadecuadas pueden provocar el golpe de mástil en baja G. Un golpe de mástil severo en vuelo suele terminar con la separación del eje del rotor principal o con el contacto de una pala con el fuselaje.

\textbf{El rotor debe volver a cargarse antes de que el cíclico lateral pueda detener cualquier alabeo.} Para recargarlo, aplique de inmediato cíclico atrás suave, evitando grandes acciones de cíclico.

Nunca intente demostrar o experimentar con maniobras de baja G, cualquiera sea su habilidad o experiencia. Incluso pilotos de prueba muy experimentados han muerto investigando la condición de baja G. Evite siempre con sumo cuidado cualquier maniobra o condición que pueda producir baja G. Los accidentes por golpe de mástil en baja G son casi siempre fatales.

Los empujes no son la única causa de baja G: la respuesta inadecuada del piloto a la turbulencia también puede introducir una peligrosa pérdida de carga del rotor (vea la SN-32).

\begin{center}\textbf{¡NUNCA HAGA UN EMPUJE EN BAJA G!}\end{center}

\subsection*{SN-13 \quad No fije objetos a los patines}
{\small\textit{Emitida: enero de 1983. Revisada: junio de 1994.}}

Los codos de los montantes del tren se han agrietado en varios helicópteros cuando el piloto intentó llevar una carga externa amarrada a los patines. El tren está optimizado para soportar grandes cargas ``hacia arriba'' y, en consecuencia, tiene muy poca resistencia en el sentido opuesto, ``hacia abajo''. Además, incluso un peso pequeño fijado al tren puede cambiar su frecuencia natural lo suficiente para provocar cargas altas por la vibración en vuelo. No intente llevar ninguna carga externa u objeto fijado al tren de aterrizaje.

\subsection*{SN-15 \quad Quedarse sin combustible puede ser fatal}
{\small\textit{Emitida: agosto de 1983. Revisada: junio de 1994.}}

Muchos pilotos subestiman la gravedad de quedarse sin combustible: equivale a una falla total y súbita del motor o de la transmisión. El piloto debe entrar de inmediato en autorrotación y prepararse para un aterrizaje forzoso (vea ``Falla de potencia'' en el capítulo~\ref{cap:emergencias}). Si no entra de inmediato en autorrotación, la RPM decaerá rápidamente, el rotor entrará en pérdida y el resultado probablemente será fatal. Se han producido accidentes graves o fatales por agotamiento del combustible. Para que no le ocurra:

\begin{pasos}
\item Nunca confíe solamente en el indicador de combustible ni en la luz de bajo combustible: estos dispositivos electromecánicos son de confiabilidad dudosa en cualquier avión o helicóptero. Anote siempre la lectura del horómetro cada vez que llene los tanques.
\item Durante el prevuelo: (a) revise visualmente el nivel en los tanques; (b) asegúrese de que las tapas estén apretadas; (c) drene una pequeña cantidad de cada tanque y del gascolator para buscar agua u otra contaminación.
\item Antes del despegue: (a) asegúrese de que la válvula de combustible esté completamente abierta; (b) verifique que el protector esté puesto en la mezcla; (c) planifique la próxima carga de modo de llegar con al menos 20 minutos de combustible.
\item En vuelo: (a) vigile continuamente el horómetro y los indicadores; si cualquiera indica poco combustible, \textbf{ATERRICE}; (b) aterrice siempre para cargar antes de que el indicador del tanque principal marque menos de 1/4; (c) \textbf{NUNCA} deje que el combustible baje tanto en vuelo que se encienda la luz de bajo combustible.
\end{pasos}

\subsection*{SN-16 \quad Los tendidos eléctricos son mortales}
{\small\textit{Emitida: abril de 1984. Revisada: junio de 1994.}}

Chocar contra cables, tirantes y otros objetos es, por lejos, la primera causa de accidentes fatales en helicópteros. Los pilotos deben estar constantemente alerta ante este peligro muy real.

\begin{itemize}
\item Busque las torres: no verá los cables a tiempo.
\item Cruce los tendidos volando directamente sobre las torres.
\item Tenga en cuenta los cables de guardia (tierra), más delgados y normalmente invisibles, que van bastante por encima de los cables más grandes y visibles.
\item Observe constantemente el terreno más alto a ambos lados de su trayectoria en busca de torres.
\item Mantenga siempre al menos 500~ft AGL, salvo en el despegue y el aterrizaje. Volando siempre sobre 500~ft AGL se elimina prácticamente la principal causa de accidentes fatales.
\end{itemize}

\subsection*{SN-17 \quad Nunca abandone el helicóptero con el motor en marcha}
{\small\textit{Emitida: noviembre de 1984. Revisada: febrero de 1989; junio de 1994.}}

Se han producido varios accidentes cuando pilotos dejaron momentáneamente sus helicópteros desatendidos con el motor en marcha y los rotores girando. El colectivo puede subir solo, aumentando el paso y el acelerador, y el helicóptero puede despegar o volcar sin control.

\textbf{Sujete los controles al embarcar pasajeros.} Es importante sujetar firmemente el cíclico y el acelerador al subir o bajar pasajeros con el motor en marcha, por si golpean inadvertidamente los controles o se deslizan sobre el acelerador abriéndolo.

\textbf{Nunca aterrice en pasto alto y seco.} El escape del motor está muy caliente y puede encender fácilmente pasto alto o matorrales. Un R22 quedó completamente destruido por el fuego tras un aterrizaje normal en pasto alto.

\subsection*{SN-18 \quad La pérdida de visibilidad puede ser fatal}
{\small\textit{Emitida: enero de 1985. Revisada: febrero de 1989; junio de 1994.}}

Volar un helicóptero con visibilidad reducida por niebla, nieve, techo bajo o incluso una noche oscura puede ser fatal. Los helicópteros tienen menos estabilidad inherente y velocidades de alabeo y cabeceo mucho mayores que los aviones. Perder las referencias visuales exteriores, aunque sea por un momento, puede provocar desorientación, acciones erróneas de los controles y un accidente sin control. Esta situación suele ocurrir cuando un piloto intenta atravesar un área parcialmente oscurecida y se da cuenta demasiado tarde de que está perdiendo visibilidad. Pierde el control al intentar un viraje para recuperarla y no logra completarlo sin referencias visuales.

¡Debe actuar \emph{antes} de perder la visibilidad! Recuerde que, a diferencia del avión, el helicóptero tiene la capacidad única de aterrizar y permitirle usar otro medio de transporte cuando hay mal tiempo, siempre que tenga el buen juicio y la voluntad necesarios para tomar la decisión correcta.

\textbf{El exceso de confianza predomina en los accidentes.} El rasgo personal más frecuente en los pilotos que sufren accidentes graves es el exceso de confianza. Los pilotos de ala fija con muchas horas que pasan a helicópteros y los propietarios privados son especialmente susceptibles. Los pilotos de avión se sienten seguros y relajados en el aire, pero aún no han desarrollado el tacto, la coordinación y la sensibilidad que exige un helicóptero. Los propietarios privados son sus propios jefes y pueden volar sin disciplina, sin reglas impuestas ni controles y críticas periódicas de un jefe de pilotos; deben depender de la autodisciplina, que a veces se olvida.

Volados de manera correcta y conservadora, los helicópteros son potencialmente las aeronaves más seguras que existen. Pero también son probablemente las menos indulgentes. Deben volarse siempre a la defensiva: el piloto debe darse un margen de seguridad mayor del que cree necesario, por si acaso.

\subsection*{SN-19 \quad Volar bajo sobre el agua es muy peligroso}
{\small\textit{Emitida: julio de 1985. Revisada: junio de 1994.}}

Muchos accidentes de helicóptero han ocurrido maniobrando bajo sobre el agua. Muchos pilotos no se dan cuenta de que pierden la percepción de profundidad sobre el agua. El agua espejada y en calma es particularmente peligrosa, pero incluso el agua picada, con su superficie siempre cambiante, interfiere con la percepción normal de profundidad y puede hacer que el piloto juzgue mal su altura.

\textbf{Mantenga 500~ft AGL siempre que sea posible y evite maniobrar sobre el agua por debajo de 200~ft AGL.}

\subsection*{SN-20 \quad Cuidado con los vuelos de demostración o de instrucción inicial}
{\small\textit{Emitida: septiembre de 1985. Revisada: junio de 1994.}}

Una proporción desmedida de accidentes fatales y no fatales ocurre en vuelos de demostración o de instrucción inicial, porque se permite que personas distintas del piloto manipulen los controles sin estar debidamente preparadas o instruidas.

Si un alumno empieza a perder el control, un instructor experimentado puede recuperarlo fácilmente siempre que el alumno no haga movimientos grandes o bruscos. Pero si el alumno se confunde por un momento y hace una acción grande y súbita en la dirección equivocada, ni siquiera el instructor más experimentado podrá recuperar el control. Los instructores suelen estar preparados para el alumno que pierde el control y no hace nada, pero rara vez para el que pierde el control y hace lo incorrecto.

Antes de permitir que alguien toque los controles, debe instruírsele a fondo sobre la extrema sensibilidad de los controles de un helicóptero liviano, e indicársele con firmeza que nunca haga un movimiento grande o súbito. Y el piloto al mando debe estar preparado para tomar instantáneamente los controles si el alumno empieza a hacer un movimiento equivocado.

\subsection*{SN-22 \quad Estado de anillo de vórtice: cómo evitarlo, reconocerlo y recuperarse}
{\small\textit{Emitida: julio de 1986. Revisada: junio de 1994; octubre de 2016; enero de 2024.}}

Un descenso vertical o una aproximación empinada, sobre todo con viento de cola, puede hacer que el rotor vuele dentro de su propio flujo descendente. A ciertas razones de descenso se forman grandes vórtices a medida que el flujo descendente recircula a través del disco rotor. Esta condición se llama estado de anillo de vórtice (VRS).

Una vez establecido el VRS, agregar potencia (subir el colectivo) puede aumentar inesperadamente la razón de descenso por el aumento del flujo que recircula por el rotor. La recuperación solo se logra sacando el disco rotor de su propio flujo descendente.

Para evitar el VRS, reduzca la razón de descenso antes de reducir la velocidad. Una buena regla es \textbf{no dejar nunca que la velocidad baje de 30 nudos hasta que la razón de descenso sea menor que 300~ft/min}.

Los signos de que se está desarrollando el VRS incluyen aumento de la vibración, menor autoridad de los controles (``controles esponjosos'') y un aumento rápido de la razón de descenso. Los pilotos deben estar siempre conscientes del viento y planificar los descensos para evitar el VRS, y estar especialmente alerta en operaciones de estacionario OGE o en aproximaciones empinadas.

La recuperación debe iniciarse en cuanto se sospeche el VRS. Es esencial reconocerlo temprano y recuperarse de inmediato sacando el rotor de su flujo descendente. No se necesitan grandes acciones de los controles y deben evitarse. Después de la recuperación, aumentar el colectivo o la velocidad ayuda a no volver a entrar en VRS.

\subsection*{SN-23 \quad Caminar hacia el rotor de cola puede ser fatal}
{\small\textit{Emitida: julio de 1986. Revisada: junio de 1994.}}

Pasajeros que no eran pilotos han muerto al caminar inadvertidamente hacia un rotor de cola girando. El piloto debe tomar todas las precauciones posibles para evitar este trágico tipo de accidente. Observe siempre estas reglas:

\begin{pasos}
\item Nunca permita que nadie se acerque al helicóptero sin escolta o sin haber sido debidamente instruido. Si es necesario, detenga el motor y los rotores antes de embarcar pasajeros.
\item Mantenga siempre la luz estroboscópica destellando mientras los rotores giren.
\item Instruya a los pasajeros para que establezcan y mantengan contacto visual con el piloto al acercarse al helicóptero (así solo se acercarán por la nariz o el costado, nunca por la cola).
\item Instruya a los pasajeros para que se alejen del helicóptero a plena vista del piloto y caminen solo alrededor de la nariz, nunca de la cola.
\item Tenga especial cuidado al aterrizar fuera de aeródromos: niños o adultos no vistos podrían acercarse al helicóptero por detrás.
\end{pasos}

\subsection*{SN-24 \quad La pérdida del rotor por baja RPM puede ser fatal}
{\small\textit{Emitida: septiembre de 1986. Revisada: junio de 1994.}}

La pérdida del rotor por baja RPM causa un porcentaje muy alto de accidentes de helicóptero, fatales y no fatales. A menudo mal entendida, no debe confundirse con la pérdida de la punta de la pala en retroceso, que ocurre solo a altas velocidades de avance cuando entra en pérdida una pequeña parte de la punta de la pala que retrocede. La pérdida de la pala en retroceso causa vibración y problemas de control, pero el rotor sigue siendo muy capaz de sustentar el peso del helicóptero.

La pérdida del rotor, en cambio, puede ocurrir a cualquier velocidad y, cuando ocurre, el rotor deja de producir la sustentación necesaria y el helicóptero literalmente cae del cielo. Afortunadamente, estos accidentes ocurren con mayor frecuencia cerca del suelo, durante el despegue o el aterrizaje, y el helicóptero cae solo cuatro o cinco pies: se destruye pero los ocupantes sobreviven. Sin embargo, la pérdida del rotor también ocurre a mayor altura, y si sucede por encima de 40 o 50~ft AGL lo más probable es que sea fatal.

La pérdida del rotor es muy parecida a la pérdida del ala de un avión a baja velocidad. A medida que el avión pierde velocidad, el ángulo de ataque del ala debe aumentar para producir la sustentación necesaria. A un ángulo crítico (unos 15°), el flujo se desprende del ala y entra en pérdida, con una súbita pérdida de sustentación y un gran aumento de la resistencia. El piloto de avión se recupera bajando la nariz para reducir el ángulo de ataque y agregando potencia para recuperar la velocidad.

En el helicóptero ocurre lo mismo, pero por baja RPM del rotor en lugar de baja velocidad. A medida que baja la RPM, el ángulo de ataque de las palas debe aumentar para generar la sustentación necesaria. Aunque el piloto no suba el colectivo, el helicóptero empezará a descender hasta que el movimiento ascendente del aire hacia el rotor proporcione el aumento de ángulo de ataque necesario. Igual que con el ala, el perfil de la pala entrará en pérdida a un ángulo crítico, con una súbita pérdida de sustentación y un gran aumento de la resistencia. La mayor resistencia actúa como un enorme freno rotor que hace caer rápidamente la RPM, agravando la pérdida. Cuando el helicóptero empieza a caer, el aire que sube sigue aumentando el ángulo de ataque de las palas, que giran lentamente, y la recuperación se vuelve prácticamente imposible, incluso con el colectivo completamente abajo.

La pérdida del rotor no es simétrica, porque cualquier velocidad de avance produce un flujo mayor en la pala que avanza que en la que retrocede. La pala que retrocede entra primero en pérdida y baja al ir hacia atrás, mientras la que avanza sigue subiendo al ir hacia adelante. La pala trasera baja y la delantera alta se traducen en una rápida inclinación hacia atrás del disco, a veces llamada ``soplado del rotor'' (\emph{rotor blow-back}). Además, cuando el helicóptero empieza a caer, el aire que sube bajo las superficies de cola tiende a bajar la nariz. Estos dos efectos, sumados al cíclico atrás que aplica el piloto para que la nariz no caiga, frecuentemente hacen que las palas se soplen hacia atrás y corten el cono de cola mientras el helicóptero cae en pérdida. Por la magnitud de las fuerzas y la flexibilidad de las palas, los topes de balanceo no impiden el corte del cono de cola. Ese corte, sin embargo, es anecdótico: la aeronave y sus ocupantes ya están condenados por el rotor en pérdida antes de que ocurra.

\subsection*{SN-25 \quad Hielo en el carburador}
{\small\textit{Emitida: diciembre de 1986. Revisada: julio de 2012.}}

Se han atribuido accidentes evitables a la detención del motor por hielo en el carburador. Usados correctamente, la calefacción del carburador y el asistente de calefacción del R22 y el R44 evitan el hielo.

Las caídas de presión y la evaporación del combustible dentro del carburador producen un enfriamiento importante, por lo que puede formarse hielo con OAT de hasta 30~°C (86~°F). Incluso con aire generalmente seco, condiciones locales como una masa de agua cercana pueden ser propicias. Ante la duda, suponga que las condiciones son propicias para el hielo y aplique calefacción según sea necesario.

En el R22 y el R44 puede ser necesaria la calefacción durante el despegue. A diferencia de los aviones, que despegan con el acelerador a fondo, los helicópteros despegan con la potencia necesaria, lo que los hace vulnerables al hielo en el carburador. Use también calefacción total durante la prueba de motor para precalentar el sistema de inducción.

En los helicópteros con asistente de calefacción, el pomo debe dejarse destrabado salvo que sea obvio que las condiciones no son propicias.

La calefacción reduce la potencia del motor para una presión de admisión dada. Con calefacción total se necesitan aproximadamente 1,5~inHg adicionales de MAP para generar la potencia máxima continua (MCP) o de despegue (TOP). Esa MAP adicional no sobrecarga el motor ni el helicóptero, porque se siguen respetando los límites de potencia. Como el motor está reducido de potencia, produce la TOP a menor altitud incluso con calefacción total. Sin embargo, evite usar más calefacción de la necesaria en altura, porque el motor puede llegar al acelerador a fondo con menos de MCP o TOP.

\subsection*{SN-26 \quad Vuelo nocturno con mal tiempo puede ser mortal}
{\small\textit{Emitida: enero de 1987. Revisada: junio de 1994.}}

Muchos accidentes fatales han ocurrido de noche cuando el piloto intentó volar en tiempo marginal después del anochecer. La tasa de accidentes fatales de noche es muchas veces mayor que de día.

De noche el piloto no puede ver los cables ni la base de las nubes, ni jirones bajos o niebla. Aunque los vea, no puede juzgar su altura porque no hay horizonte de referencia. No se da cuenta de que están allí hasta que vuela dentro de ellos y pierde de golpe sus referencias visuales exteriores y la capacidad de controlar la actitud del helicóptero. Como los helicópteros no son inherentemente estables y tienen velocidades de alabeo muy altas, la aeronave queda rápidamente fuera de control y el resultado es un impacto a alta velocidad, normalmente fatal.

Asegúrese de \textbf{NUNCA} volar de noche salvo con tiempo despejado, techo ilimitado o muy alto y abundante luz celeste o terrestre como referencia.

\subsection*{SN-27 \quad Los cortes de acelerador por sorpresa pueden ser mortales}
{\small\textit{Emitida: diciembre de 1987. Revisada: junio de 1994.}}

Muchos instructores no saben dar a un alumno una falla de potencia simulada de manera segura. Quizás aprendieron a responder ellos mismos a un corte de acelerador, pero no a preparar al alumno ni a manejar una reacción inesperada. El alumno puede congelarse sobre los controles, pisar el pedal equivocado, subir el colectivo en lugar de bajarlo o simplemente no hacer nada. El instructor debe estar preparado para cualquier reacción inesperada.

Antes de dar una falla de potencia simulada, prepare cuidadosamente al alumno y asegúrese de haber volado juntos lo suficiente para establecer el entendimiento y la comunicación críticos entre instructor y alumno. Repitan el ejercicio varias veces hasta que las reacciones del alumno sean correctas y predecibles. Nunca sorprenda de verdad al alumno: dígale unos minutos antes que le dará una falla de potencia simulada y, al cerrar el acelerador, anuncie en voz alta ``falla de potencia''. La presión de admisión debe ser menor que 21~inHg y el acelerador debe cerrarse suavemente, nunca ``cortarse''. Acompañe todos los controles y tense los músculos de la pierna derecha para impedir que el alumno pise el pedal equivocado si se confunde. Y suponga siempre que tendrá que completar usted mismo la entrada en autorrotación: nunca espere a ver qué hace el alumno. Planifique iniciar la recuperación dentro de un segundo, cualquiera sea la reacción del alumno.

Ha habido casos en que el motor se detuvo durante fallas simuladas. Como precaución, haga siempre la falla simulada dentro de la distancia de planeo de un área abierta y lisa donde esté seguro de poder completar una autorrotación con contacto si fuera necesario. Además, nunca practique fallas de potencia simuladas antes de que el motor esté completamente caliente: espere a llevar al menos 15 a 20 minutos de vuelo.

\subsection*{SN-28 \quad Escuche el aviso de una falla inminente de rodamiento}
{\small\textit{Emitida: julio de 1988. Revisada: julio de 2012.}}

La falla inminente de un rodamiento de bolas o de rodillos suele estar precedida por un aumento notorio del ruido, que normalmente comienza varias horas antes de que el rodamiento falle o de que aumente su temperatura. Para detectarla, el piloto debe descubrirse un oído y escuchar el sonido del sistema de transmisión durante el arranque y la detención. Una vez familiarizado con el sonido normal, debería poder detectar el ruido de un rodamiento en falla: un fuerte zumbido, retumbo, gruñido o sonido de sirena. Al oír un ruido inusual, el piloto debe dejar de inmediato la aeronave en tierra y hacer inspeccionar a fondo los rodamientos por un mecánico calificado. La falla de un rodamiento en vuelo podría causar un accidente grave.

No confíe en los Telatemps para detectar una falla inminente de rodamiento. Un rodamiento en falla puede no calentarse lo suficiente para ennegrecer los Telatemps hasta que empieza a desintegrarse, lo que puede ocurrir solo segundos antes de la falla completa.

\textbf{Advertencia de la luz del embrague.} Es normal que la luz del embrague se encienda ocasionalmente en vuelo por poco tiempo (unos 3 a 6 segundos) para retensar las correas. Si la luz parpadea o no se apaga dentro de 10 segundos, puede indicar una falla de correa o de rodamiento. Ante una indicación anormal, saque el disyuntor del embrague y reduzca la potencia. Elija un lugar seguro y haga un aterrizaje precautorio para revisar el sistema de transmisión. Si hay otros síntomas de falla de la transmisión (olor a goma caliente, ruido o vibración), aterrice de inmediato. Si se separan las agujas del tacómetro, entre en autorrotación.

Después de aterrizar, detenga el motor y revise que las correas estén en sus ranuras y sin daños. Revise los sellos de los rodamientos superior e inferior del actuador y las lecturas de los Telatemps. Si encuentra problemas en la transmisión, haga inspeccionar la aeronave por un mecánico antes de volar de nuevo.

\subsection*{SN-29 \quad Pilotos de avión: alto riesgo al volar helicópteros}
{\small\textit{Emitida: marzo de 1993. Revisada: junio de 1994.}}

Ha habido varios accidentes fatales de pilotos con muchas horas en avión pero experiencia limitada en helicóptero.

Las reacciones arraigadas de un piloto de avión experimentado pueden ser mortales en un helicóptero. Puede volarlo bien en maniobras normales y condiciones ordinarias, cuando tiene tiempo para pensar la respuesta correcta. Pero si debe reaccionar de golpe ante algo inesperado, puede volver a sus reacciones de avión y cometer un error fatal: sus manos y pies se mueven por puro reflejo, sin pensamiento consciente, y esos reflejos bien pueden basarse en su mayor experiencia, la de volar aviones.

Por ejemplo, en un avión la reacción ante una bocina de aviso (pérdida) sería llevar de inmediato la palanca adelante y agregar potencia. En un helicóptero, llevar el cíclico adelante al oír una bocina (baja RPM) bajaría aún más la RPM y podría causar la pérdida del rotor, sobre todo si además ``agrega potencia'' (sube el colectivo). En menos de un segundo podría hacer entrar el rotor en pérdida y el helicóptero caería del cielo.

Otro ejemplo es la reacción para hacer descender la aeronave. Si el piloto de helicóptero debe descender de golpe para evitar un ave u otra aeronave, baja rápidamente el colectivo moviendo muy poco el cíclico. En la misma situación, el piloto de avión empujaría la palanca adelante para picar. Un movimiento rápido del cíclico adelante en estas condiciones produciría una condición de baja G que podría causar golpe de mástil, con la separación del eje del rotor o el golpe de una pala contra el fuselaje. Algo similar ocurre al terminar un ascenso después de un tirón: el piloto de avión lo hace con palanca adelante; el de helicóptero debe usar el colectivo o una aplicación muy gradual y suave del cíclico adelante.

Para sobrevivir en el helicóptero, el piloto de avión experimentado debe dedicar mucho tiempo y esfuerzo a desarrollar reacciones seguras de helicóptero, que deben ser más fuertes y prevalecer sobre las de avión porque en el helicóptero todo ocurre más rápido. No hay tiempo para darse cuenta del error, pensarlo y corregirlo: es demasiado tarde, el rotor ya entró en pérdida o una pala ya golpeó la célula y no hay posibilidad de recuperación. Para desarrollar reacciones seguras, debe practicar cada procedimiento una y otra vez con un instructor competente hasta que sus manos y pies siempre hagan el movimiento correcto sin pensamiento consciente. Y, \textbf{SOBRE TODO, NUNCA DEBE EMPUJAR BRUSCAMENTE EL CÍCLICO HACIA ADELANTE}. Vea también las SN-11 y SN-24.

\subsection*{SN-30 \quad Los objetos sueltos pueden ser fatales}
{\small\textit{Emitida: junio de 1994. Revisada: abril de 2009.}}

Se han producido accidentes fatales por objetos sueltos que salieron de la cabina y golpearon el rotor de cola. Cualquier objeto que golpee el rotor de cola puede provocar la falla de una pala. La pérdida o el daño de una pala del rotor de cola puede causar un desequilibrio severo capaz de separar la caja de cola o todo el conjunto de cola del cono de cola, con un accidente catastrófico. También se han producido accidentes por tapas de combustible, aves y otros objetos que golpearon el rotor de cola. Antes de cada vuelo:

\begin{pasos}
\item Camine alrededor de toda la aeronave revisando que las tapas de combustible estén firmes y el estado del rotor de cola. Asegúrese de que no haya objetos sueltos ni basura cerca del helicóptero. Verifique que todos los pasadores de bisagra de las puertas tengan sus anillos o pasadores de chaveta.
\item Guarde o asegure todos los objetos sueltos de la cabina. Incluso con puertas, objetos como cartas pueden ser succionados por una ventilación.
\item Instruya a los pasajeros sobre el peligro de que un objeto golpee el rotor de cola. Adviértales que nunca arrojen nada desde el helicóptero ni pongan objetos cerca de las ventilaciones, de donde podrían ser succionados.
\item Trabe firmemente todas las puertas.
\item Nunca vuele con la puerta izquierda retirada. (Retire solo la derecha para ventilar.)
\end{pasos}

\subsection*{SN-31 \quad El gobernador puede enmascarar el hielo en el carburador}
{\small\textit{Emitida: diciembre de 1996.}}

Con el gobernador conectado, el hielo en el carburador no se manifestará como pérdida de RPM ni de presión de admisión: el gobernador ajustará automáticamente el acelerador para mantener constante la RPM, lo que también mantiene constante la presión de admisión. Ante la duda, aplique la calefacción necesaria para mantener la CAT fuera del arco amarillo en estacionario, ascenso o crucero, y aplique calefacción total cuando la presión de admisión esté por debajo de 18~inHg. Recuerde también que, si se usa el asistente de calefacción, este reducirá la calefacción al elevarse a estacionario y el control puede requerir reajuste en vuelo.

\subsection*{SN-32 \quad Vientos fuertes o turbulencia}
{\small\textit{Emitida: marzo de 1998. Revisada: junio de 2020; julio de 2025.}}

Debe evitarse volar con vientos fuertes o turbulencia. Una aplicación inadecuada de los controles en respuesta a la turbulencia o a vientos fuertes puede provocar la pérdida de control o una peligrosa condición de baja G. El piloto debe conocer las condiciones asociadas a la turbulencia: actividad convectiva, vientos fuertes en superficie sobre terreno de elevación variable, ondas de montaña e inversiones fuertes. Evite volar sobre terreno donde la turbulencia probablemente sea severa. Si se espera o encuentra turbulencia moderada o mayor, se recomienda:

\begin{itemize}
\item Reducir la potencia y usar una velocidad de crucero menor que la normal (60--70 nudos). El golpe de mástil en baja G es mucho menos probable a menor velocidad.
\item Apretar el cinturón y apoyar el antebrazo derecho en la pierna derecha para minimizar las acciones involuntarias del cíclico. Algunos pilotos pueden optar por aplicar un poco de fricción al cíclico para minimizarlas aún más.
\item Desconectar los modos superiores del piloto automático, como ALT, VS, IAS, HDG y NAV.
\item No sobrecontrolar. Dejar que la aeronave se mueva con la turbulencia y luego restablecer la actitud con acciones suaves y delicadas. Son de esperar desviaciones momentáneas de velocidad, rumbo, altitud y RPM.
\item Si se encuentra una condición de baja G, volver a cargar el rotor con un movimiento suave del cíclico hacia atrás antes de corregir cualquier desviación de alabeo.
\item Si corresponde, ajustar el rumbo para evitar terreno que pueda causar turbulencia, como el lado de sotavento de cerros, cordones o edificios altos.
\end{itemize}

El helicóptero es más susceptible a la turbulencia con poco peso. Reduzca la velocidad y tenga cuidado al volar solo o con poca carga. No vuele si se prevén condiciones extremas. Vea la SN-11 para más información sobre la baja G.

\subsection*{SN-33 \quad Holgura de las correas de transmisión}
{\small\textit{Emitida: marzo de 1998. Revisada: julio de 2013; julio de 2021.}}

Los pilotos deben asegurarse de que las correas del R22 y del R44 no tengan holgura excesiva durante el arranque. Las correas demasiado sueltas pueden saltar de las ranuras de las poleas al tensarse, y el piloto no tiene forma de observarlo. Por eso estos ítems de la lista son muy importantes:

\begin{pasos}
\item Según la lista de prevuelo, con el embrague desacoplado, desvíe las correas hacia el interior justo sobre la carcasa del ventilador, en el lado derecho de la aeronave. Verifique una desviación máxima de 1,5~in (4~cm). Si se desvían más, haga ajustar el actuador de tensión de correas antes de volar.
\item Según la lista de prueba de motor, verifique que el rotor gire dentro de 5 segundos de acoplar el interruptor del embrague. Si tarda más, detenga el motor y haga ajustar el actuador antes de volar.
\end{pasos}

\subsection*{SN-34 \quad Vuelos de levantamiento aéreo y fotográficos: riesgo muy alto}
{\small\textit{Emitida: marzo de 1999. Revisada: abril de 2009.}}

Existe la idea errónea de que los vuelos de levantamiento aéreo y fotográficos pueden hacerlos con seguridad pilotos de poca experiencia. No es así: ha habido numerosos accidentes fatales en estos vuelos, varios con helicópteros Robinson.

A menudo, para complacer al observador o al fotógrafo, un piloto inexperto reduce la velocidad a menos de 30~KIAS e intenta maniobrar para el mejor ángulo de visión. Mientras maniobra puede perder de vista la velocidad y el viento. El helicóptero puede perder rápidamente la sustentación traslacional y empezar a hundirse. Un piloto inexperto puede subir el colectivo para detener el descenso, lo que puede bajar la RPM, reduciendo la potencia disponible y causando una razón de descenso aún mayor y más pérdida de RPM. Abrir el acelerador aumentará el torque del rotor pero no la potencia disponible, por la baja RPM. Como el empuje del rotor de cola es proporcional al cuadrado de la RPM, si la RPM cae por debajo de 80\,\% se pierde casi la mitad del empuje del rotor de cola y el helicóptero gira con la nariz a la derecha. De pronto, la caída de RPM también hace entrar en pérdida al rotor principal y el helicóptero cae rápidamente mientras sigue girando. El impacto suele ser fatal.

Los vuelos de levantamiento aéreo y fotográficos solo deben hacerlos pilotos bien entrenados y experimentados que:

\begin{pasos}
\item Tengan al menos 500 horas como piloto al mando en helicópteros y más de 100 horas en el modelo que vuelan.
\item Tengan amplio entrenamiento en técnicas de recuperación de baja RPM y de asentamiento con potencia (\emph{settling with power}).
\item Estén dispuestos a decir que no al observador o fotógrafo y vuelen solo a velocidades, alturas y ángulos de viento seguros que permitan buenas rutas de escape.
\end{pasos}

Vea también la SN-24.

\subsection*{SN-35 \quad Vuelo cerca de torres de radiodifusión}
{\small\textit{Emitida: abril de 1999. Revisada: febrero de 2016.}}

Las torres de radiodifusión de alta potencia (típicamente de 50~kW o más) pueden generar interferencia electromagnética (EMI) con los sistemas eléctricos y equipos electrónicos del helicóptero. Los primeros indicios de un campo de radio de alta potencia pueden ser una considerable apertura de \emph{squelch} o estática en el intercomunicador. Efectos más severos pueden incluir el encendido aleatorio de luces de advertencia o un funcionamiento errático del gobernador y del tacómetro. En el peor caso, el gobernador puede intentar abrir o cerrar completamente el acelerador; si el piloto no está listo para contrarrestarlo, puede producirse una baja RPM o una sobrevelocidad.

La intensidad del campo disminuye rápidamente con la distancia a la antena. Para minimizar la probabilidad de EMI, \textbf{no vuele a menos de media milla de torres de radiodifusión de alta potencia}. Si encuentra EMI inadvertidamente:

\begin{itemize}
\item No se distraiga intentando ajustar la radio o el sistema de audio.
\item Mantenga una mano en la empuñadura del acelerador y sienta su funcionamiento normal. Si es anormal, apriete la empuñadura para sobreponerse al gobernador y controle el acelerador manualmente. Si lo desea, desconecte el gobernador hasta salir del área de EMI.
\item Vigile atentamente el tacómetro, los instrumentos de motor y las luces de advertencia, sin reaccionar en exceso ante una indicación falsa.
\item Después de un vuelo con EMI, haga evaluar el sistema eléctrico y el equipo electrónico por personal de mantenimiento calificado.
\end{itemize}

La ubicación y altura de las torres figuran en las cartas aeronáuticas, pero no su potencia ni su frecuencia. Aunque todas deben rodearse a distancia segura, las de la banda HF/onda corta son las más propensas a causar EMI. En el sitio web de Robinson (\texttt{www.robinsonheli.com/sn35antennas}) hay una lista de esas torres en los EE.\,UU., compilada de la base de datos de la FCC; los organismos correspondientes de otros países pueden tener bases de datos similares.

\subsection*{SN-36 \quad Sobrevelocidades durante el despegue}
{\small\textit{Emitida: noviembre de 2000.}}

Helicópteros han sufrido daños graves por sobrevelocidades de RPM durante el despegue. Las sobrevelocidades causaron una vibración del eje del rotor de cola que provocó la falla inmediata del eje y del cono de cola. En todo el rango normal de RPM, la vibración del eje la controla el rodamiento amortiguador, pero este no es eficaz por encima de 120\,\% de RPM.

La correlación mecánica puede causar una sobrevelocidad en el despegue si la RPM se lleva a valores normales de vuelo y se sube el colectivo antes de conectar el gobernador. También puede ocurrir si se aprieta demasiado el acelerador durante el despegue, sobreponiéndose al gobernador. Los pilotos inexpertos, más propensos a estar nerviosos o distraídos, son particularmente susceptibles. Para evitarlas:

\begin{pasos}
\item Confirme siempre el gobernador conectado antes de aumentar la RPM sobre 80\,\%.
\item Verifique que el gobernador estabilice la RPM del motor cerca del extremo superior del arco verde.
\item Mantenga el acelerador tomado sin apretar, dejando que el gobernador controle la RPM.
\end{pasos}

\subsection*{SN-37 \quad Exceder las limitaciones aprobadas puede ser fatal}
{\small\textit{Emitida: diciembre de 2001.}}

Muchos pilotos no entienden la fatiga del metal. Cada vez que un componente metálico se carga por encima de su límite de fatiga, se produce un daño oculto dentro del metal. No existe método de inspección que detecte este daño invisible. El primer indicio será una grieta microscópica, a menudo oculta, que crecerá con cada repetición del esfuerzo crítico hasta que la pieza se rompa de golpe. El crecimiento es muy rápido en las piezas de la transmisión, por las cargas torsionales de alta frecuencia, y también en los componentes del sistema rotor, por la alta fuerza centrífuga sobre las palas y el cubo. Con cada vuelta de un eje o una pala sobrecargados se producen ciclos de fatiga dañinos.

Si un piloto excede los límites de potencia o velocidad algunas veces sin que nada falle, puede creer equivocadamente que puede operar con seguridad a esas cargas. No es así: cada segundo que se exceden las limitaciones se producen más ciclos de esfuerzo y puede acumularse más daño. Tarde o temprano empezará una grieta de fatiga que crecerá hasta una falla súbita. Con suerte, la pieza habrá alcanzado su vida aprobada y se habrá reemplazado antes; si no, probablemente habrá un accidente grave o fatal.

\begin{advertencia}
\begin{pasos}
\item Opere siempre bastante por debajo de la \textsc{v}\textsubscript{NE} aprobada, especialmente con viento turbulento.
\item No opere el motor por encima de los límites de presión de admisión de los letreros.
\item No cargue la aeronave por encima del peso bruto aprobado.
\item Las condiciones más dañinas se dan al volar o maniobrar a alta velocidad combinada con alta potencia.
\end{pasos}
\end{advertencia}

\subsection*{SN-38 \quad Las autorrotaciones de práctica causan muchos accidentes de instrucción}
{\small\textit{Emitida: julio de 2003. Revisada: julio de 2019.}}

Las autorrotaciones de práctica siguen siendo la primera causa de accidentes de instrucción en helicópteros, y muchos ocurren porque el instructor no toma una acción decidida. Los instructores deben hacer siempre el siguiente \textbf{control de decisión a 100~ft}:

\begin{pasos}
\item RPM del rotor aproximadamente 100\,\%.
\item Velocidad estabilizada entre 60 y 70~KIAS.
\item Razón de descenso estabilizada, por lo general por debajo de 1500~ft/min.
\item Virajes (si los hay) completados.
\end{pasos}

Antes de descender bajo 100~ft AGL, el instructor debe decidir tomar los controles y hacer una recuperación de potencia inmediata si cualquiera de estos parámetros no es correcto. No intente salvar la situación dando indicaciones al alumno o corrigiendo por debajo de 100~ft AGL. Con altitudes de densidad sobre 4000~ft, eleve el control de decisión a 200~ft AGL o más. Si decide continuar, el instructor debe anunciar ``control de 100~ft completado, continúe''.

Un alto porcentaje de los accidentes de instrucción ocurre después de muchas autorrotaciones consecutivas. Para mantener la concentración del instructor y minimizar la fatiga del alumno, limite la práctica a no más de 3 o 4 autorrotaciones consecutivas.

El propósito de la autorrotación de práctica es enseñar el control y la maniobra durante el planeo, el uso correcto de la recogida y la recuperación de potencia (o el aterrizaje). Si bien la práctica de fallas de potencia simuladas y aterrizajes forzosos puede implicar una reducción de potencia más brusca, las entradas en autorrotación de práctica deben hacerse de manera suave y lenta, según la \emph{Robinson Flight Training Guide}.

\subsection*{SN-39 \quad Una vibración inusual puede indicar una grieta en una pala del rotor principal}
{\small\textit{Emitida: julio de 2003.}}

Puede evitarse una falla catastrófica por fatiga de una pala si pilotos y mecánicos están atentos a los primeros indicios de una grieta. Aunque la grieta puede ser interna y no visible, probablemente causará un aumento significativo de la vibración del rotor antes de la falla final. Si un rotor queda suave después de balancearlo pero se desbalancea de nuevo en pocos vuelos, debe considerarse sospechoso: haga examinar a fondo el sistema rotor por un mecánico calificado antes de volar de nuevo.

Si la vibración del rotor principal aumenta rápidamente o se vuelve severa en vuelo, haga un aterrizaje seguro inmediato. No intente continuar el vuelo hasta un destino conveniente.

\subsection*{SN-40 \quad Incendios posteriores al accidente}
{\small\textit{Emitida: julio de 2006. Revisada: mayo de 2013.}}

Ha habido varios casos en que los ocupantes de helicópteros o avionetas sufrieron quemaduras graves por un incendio después de un accidente. La ropa de vuelo ignífuga de Nomex reduce la probabilidad de quemaduras graves. Organizaciones militares, policiales, médicas y otras suelen exigirla en todo vuelo. Los pilotos deben considerar los beneficios de la ropa ignífuga e instruir o equipar a todos los ocupantes en consecuencia.

\subsection*{SN-41 \quad Distracciones del piloto}
{\small\textit{Emitida: mayo de 2013. Revisada: julio de 2021; julio de 2025.}}

Las distracciones del piloto han causado o contribuido a numerosos accidentes fatales. Ejemplos son el uso de teléfonos móviles y dispositivos electrónicos portátiles, los eventos inesperados, comer o beber y atender a los pasajeros.

Las distracciones pueden ser visuales, físicas o mentales: todas desvían la atención del piloto de su deber principal, volar el helicóptero. Numerosos estudios muestran que las distracciones al conducir equivalen a conducir bajo influencia de sustancias, lo que llevó a leyes que lo prohíben. Operar una aeronave distraído es igualmente peligroso.

Es su responsabilidad como piloto mantener la concentración en la operación del helicóptero, sin realizar actividades que lo distraigan innecesariamente. El uso del piloto automático nunca libera al piloto de su responsabilidad de vigilar y controlar de cerca la aeronave: no caiga en la complacencia ni lo vea como una oportunidad para actividades no esenciales. Estas pautas ayudarán a su seguridad y la de sus pasajeros:

\begin{itemize}
\item Mantenga las manos disponibles para tomar el control de la aeronave en todo momento.
\item Posponga las actividades de menor prioridad hasta que sea seguro o hasta después de aterrizar.
\item Mantenga siempre su exploración visual hacia afuera y de los instrumentos de vuelo.
\item Mantenga la conciencia situacional y el compromiso mental con la operación del helicóptero.
\item Toda programación de aviónica que tome más de unos segundos debe hacerse en tierra.
\item Si una puerta se abre en vuelo, aterrice siempre antes de intentar cerrarla.
\item En estacionario, mantenga ambas manos en los controles. Si debe sintonizar una radio u otra tarea, primero aterrice y reduzca el paso colectivo.
\end{itemize}

\subsection*{SN-42 \quad Guiñada imprevista}
{\small\textit{Emitida: mayo de 2013. Revisada: julio de 2019.}}

Si el piloto no aplica los pedales correctos en respuesta a vientos fuertes o racheados durante el estacionario o el vuelo a baja velocidad, puede producirse una guiñada imprevista. Algunos pilotos la atribuyen erróneamente a la pérdida de efectividad del rotor de cola (LTE), lo que implicaría que el rotor de cola entró en pérdida o no pudo dar empuje suficiente. Los rotores de cola de los helicópteros Robinson están diseñados con más autoridad que muchos otros y es poco probable que sufran LTE.

Para evitar la guiñada imprevista, el piloto debe conocer las condiciones (un viento cruzado por la izquierda, por ejemplo) que pueden exigir acciones de pedal grandes o rápidas. Practicar giros en estacionario lentos y a velocidad constante ayuda a mantener la competencia en el control de guiñada, y el entrenamiento en estacionario con un instructor calificado en distintas condiciones de viento también puede ayudar.

Note que el empuje de cualquier rotor de cola disminuye significativamente al bajar la RPM. Una RPM baja combinada con un torque alto, como ocurre con exceso de paso, puede producir una guiñada a la derecha incontrolable (vea también la SN-34).

\subsection*{SN-43 \quad Extreme las precauciones en los vuelos posteriores al mantenimiento}
{\small\textit{Emitida: enero de 2015.}}

Varios accidentes fatales ocurrieron en vuelos inmediatamente posteriores a un mantenimiento. En varios casos la causa fue un rearmado incorrecto o incompleto del helicóptero, y el error habría sido detectable en una inspección prevuelo cuidadosa.

Incluso el mejor personal de mantenimiento puede distraerse y equivocarse. Los pilotos deben hacer un prevuelo especialmente minucioso después de un mantenimiento. Si es posible, hable con los técnicos que hicieron el trabajo, averigüe exactamente qué se hizo y preste especial atención a esas áreas. El personal profesional apreciará el compromiso del piloto con la seguridad y agradecerá una revisión adicional de su trabajo.

Todo trabajo en el sistema de controles de vuelo merece especial atención, porque una desconexión de controles es casi siempre catastrófica. Durante los trabajos de \emph{tracking} y balanceo, suba siempre a la cabeza del rotor para inspeccionar de cerca los sujetadores de las varillas de paso y de los tubos de control después de cada ajuste. Nunca apure ni omita pasos del prevuelo.

\subsection*{SN-44 \quad Transporte de pasajeros}
{\small\textit{Emitida: febrero de 2017. Revisada: junio de 2017.}}

Llevar pasajeros es una responsabilidad adicional para el piloto al mando. Los pasajeros han depositado su confianza enteramente en el piloto y deben ser informados de los riesgos asociados al vuelo, entre ellos la experiencia del piloto, la capacidad de la aeronave y consideraciones operacionales como el vuelo sobre el agua o nocturno. Llevar un pasajero, en sí mismo, aumenta el riesgo, porque agrega carga de trabajo y distracciones.

\textbf{Instrucción a los pasajeros.} Haga siempre una instrucción de seguridad antes de permitir que un pasajero suba a bordo. Debe incluir:

\begin{itemize}
\item \textbf{Riesgos:} una conversación sobre los riesgos del vuelo. Asegúrese de que los pasajeros sepan que usted cancelará o interrumpirá el vuelo con gusto si alguien se siente incómodo.
\item \textbf{Embarque y desembarque:} que los pasajeros demuestren el uso correcto de cinturones y puertas. Explique las rutas seguras de entrada y salida de emergencia (lejos del rotor de cola y a la vista del piloto). No embarque ni desembarque pasajeros con los rotores girando, salvo que se hayan establecido procedimientos (como escoltas en tierra) antes del vuelo.
\item \textbf{Controles de la aeronave:} retire siempre los controles del lado del pasajero. Asegúrese de que los pasajeros aseguren sus objetos personales. Advierta que no golpeen inadvertidamente el poste central del cíclico. Asegúrese de que sepan usar los audífonos y el intercomunicador.
\item \textbf{Expectativas del vuelo:} quizás no se pueda cumplir un horario fijo o sobrevolar ciertos lugares por meteorología, espacio aéreo o limitaciones de performance. Si los pasajeros lo saben de antemano, habrá menos presión sobre el piloto para intentar algo más allá de su comodidad o experiencia.
\item \textbf{Momentos de alta carga de trabajo:} los pasajeros siempre deben avisar al piloto de asuntos de seguridad, como tránsito, pero abstenerse de conversación innecesaria durante el despegue, el aterrizaje y las comunicaciones radiales.
\end{itemize}

\textbf{Pilotos recién habilitados.} Los datos muestran que los pilotos nuevos que llevan pasajeros tienen altas tasas de accidentes: la carga de trabajo y las distracciones que agregan los pasajeros se magnifican en quien todavía se está acostumbrando a la aeronave y al entorno. Los pilotos de helicóptero recién habilitados:

\begin{itemize}
\item No deberían llevar pasajeros hasta acumular 100 horas totales de experiencia en helicóptero, incluidas 20 horas como piloto al mando en el tipo después de obtener la habilitación.
\item Deberían volar solo de día, en buenas condiciones VFR y con poco viento.
\item En los primeros vuelos con pasajeros, deberían limitarse a un área local conocida.
\end{itemize}

\textbf{Safety Notices relacionadas:} SN-30 (no se recomienda el vuelo sin puertas con pasajeros, salvo que tengan bastante familiaridad con helicópteros) y SN-34 (las misiones fotográficas tienen riesgos muy superiores a los vuelos normales con pasajeros y requieren entrenamiento y experiencia específicos).
